\TieudeT{PHONG TỎA VẬT LÍ 11 - DAO ĐỘNG VÀ SÓNG}
\subsubsection*{PHẦN I. Thí sinh trả lời câu hỏi từ câu 1 đến câu 18. Mỗi câu hỏi thí sinh chỉ chọn một phương án}
\setcounter{ex}{0}
	\begin{ex}%[3P2Y1-1][Loai1]
		Một sóng dọc truyền trong một môi trường thì phương dao động của các phần tử môi trường
		\choice
		{là phương ngang}
		{là phương thẳng đứng}
		{\True trùng với phương truyền sóng}
		{vuông góc với phương truyền sóng}
		%	\loigiai{Một sóng dọc truyền trong một môi trường thì phương dao động của các phần tử môi trường trùng với phương truyền sóng}
	\end{ex}
	
	\begin{ex}%[3P2Y1-1][Loai1]
		
		Một sóng cơ hình sin truyền theo trục Ox. Công thức liên hệ giữa tốc độ truyền sóng v, bước sóng $\lambda$ và tần số f của sóng là
		\choice
		{$\lambda =\dfrac{f}{v}$}
		{\True $\lambda =\dfrac{v}{f}$}
		{$\lambda = 2\pi fv$}
		{$\lambda = vf$}
		\loigiai{Ta có: $\lambda =\dfrac{v}{f}$}
	\end{ex}
	\begin{ex}%[3P2Y1-2][Loai1] 
		Một sóng truyền trên mặt nước có bước sóng $\lambda = 2\ m$, khoảng cách giữa hai điểm gần nhau nhất trên cùng một phương truyền sóng dao động cùng pha nhau là
		\choice
		{$0,5\ m$}
		{$1,5\ m$}
		{\True $2\ m$}
		{$1\ m$}
		\loigiai{
			Khoảng cách gần nhau nhất trên cùng phương truyền sóng dao động cùng pha là $\lambda =2$ m.
		}
	\end{ex}
	\begin{ex}%[3P2Y2-1][Loai1]
		Hai điểm P và Q cùng nằm trên một phương truyền sóng và cách nhau một khoảng bằng $\dfrac{\lambda}{2}$ thì
		\choice
		{khi P có vận tốc cực đại thì Q ở li độ cực đại}
		{khi P có li độ cực đại, thì Q cũng có li độ cực đại}
		{\True li độ dao động của P và Q luôn luôn bằng nhau về độ lớn nhưng ngược dấu}
		{khi P đi qua vị trí cân bằng thì Q ở biên}
		\loigiai{
			P và Q dao động ngược pha $\Rightarrow \frac{u_P}{u_Q}=-1$.
		}
	\end{ex}
	%\loai{Thay số tính khoảng cách}
	\begin{ex}%[3P2B2-1][Loai1] 
		Sóng truyền từ M đến N dọc theo phương truyền sóng với bước sóng $\lambda=120\ cm$. Biết rằng sóng tại N trễ pha hơn sóng tại M là $\dfrac{\pi}{3}\ rad$. Khoảng cách $d=MN$ là
		\choice
		{$d=15\ cm$}
		{$d=24\ cm$}
		{$d=30\ cm$}
		{\True $d=20\ cm$}
		\loigiai{
			Ta có:	$\varphi = 2\pi.\frac{d}{\lambda} \Rightarrow \frac{d}{\lambda} = \frac{1}{6}\Rightarrow d = 20\ cm$.
		}
	\end{ex}
	\begin{ex}%[3P2B3-1][Loai1] 
		Một nguồn phát sóng dao động theo phương trình $u=u_0\cos10\pi t$ cm với t tính bằng giây, bước sóng là $\lambda$. Trong khoảng thời gian $2\ s$, sóng này truyền đi được quãng đường bằng
		\choice
		{$15\lambda$}
		{$5\lambda$}
		{\True $10\lambda$}
		{$20\lambda$}
		\loigiai{
			\begin{itemize}
				\item $T = \dfrac{2\pi}{\omega} = \dfrac{2\pi}{10\pi} = 0,2$ s $\Rightarrow t = 2\ s = 10T$.
				\item Trong 1T sóng truyền được quãng đường $\lambda$.
				\item 10T sóng truyền được quãng đường $S=10\lambda$.
			\end{itemize}
		}
	\end{ex}
	
	\begin{ex}%[3P2K2-2][Loai1]
		Một sóng ngang hình sin lan truyền trên trục Ox với biên độ không đổi và tần số bằng $4\ Hz$, tốc độ truyền sóng bằng $4,8\ m/s$. Trên $Ox$ có hai phần tử môi trường tại P và Q luôn chuyển động cùng tốc độ nhưng ngược chiều nhau. Trong khoảng giữa PQ có hai điểm mà phần tử môi trường tại đó dao động cùng pha với phần tử môi trường tại P. Khoảng cách PQ bằng
		\choice
		{$8\ m$}
		{$2\ m$}
		{\True $3\ m$}
		{$4\ m$}
		\loigiai{
			\begin{itemize}
				\item Ta có: $\lambda = \dfrac{v}{f} = \dfrac{4,8}{4} = 1,2\ m$\\ 
				\item Lại có $v_P = -v_Q$ hay hai phần tử P và Q dao động ngược pha với nhau
				$\Rightarrow PQ = k\lambda + \dfrac{\lambda}{2}$.
				\item Mà giữa P và Q có 2 điểm mà phần tử môi trường tại đó dao động cùng pha với phần tử môi trường tại P nên $PQ = 2\lambda + \dfrac{\lambda}{2} = 2,5\lambda$
				$\Rightarrow PQ = 2,5\lambda = 2,5.1,2 = 3$ m.
			\end{itemize}
		}
	\end{ex}
	%\loai{Thời điểm thứ n liên quan gia tốc}
	
	\begin{ex}%[3P5-4.1G][Loai9]
		Trong thí nghiệm Young về giao thoa ánh sáng, thực hiện đồng thời với hai ánh sáng đơn sắc. Khoảng vân giao thoa trên màn lần lượt là $1,2\ mm$ và $1,8\ mm$. Trên màn quan sát, gọi M, N là hai điểm ở cùng một phía so với vân trung tâm và cách vân trung tâm lần lượt là $6\ mm$ và $20\ mm$. Trên đoạn MN, quan sát được bao nhiêu vạch sáng?
		\choice
		{$19$}
		{\True $16$}
		{$20$}
		{$18$}
		\loigiai{ 
			\begin{itemize}
				\item Số vị trí vân sáng trùng nhau trên MN:\\		
				$\left\{\begin{aligned}& x=k_1i_1=k_2i_2=k_1.1,2=k_2.1,8\ mm\Rightarrow \dfrac{k_1}{k_2}=\dfrac{3}{2}\Rightarrow \left\{\begin{aligned}& k_1=3n \\
						& k_2=2n 
					\end{aligned}\right. \\
					& x=3n.1,2\ mm=3,6n\ mm\Rightarrow \left\{\begin{aligned}& 6\le x\le 20\Rightarrow 1,7\le n\le 5,6 \\
						& \Rightarrow n=2,\ 3,\ 4,\ 5:\text{ số vị trí trùng: 4} 
					\end{aligned}\right. 
				\end{aligned}\right.$
				\item Số vân sáng của hệ 1 và hệ 2 trên MN lần lượt được xác định như sau:\\
				$\left\{\begin{aligned}& 6\le x=k_1i_1=k_1.1,2\le 20\Rightarrow 5\le k_1\le 16,7\Rightarrow k_1=5,...,16:\text{ số giá trị } k_1 \text{ là 12} \\
					& 6\le x=k_2i_2=k_2.1,8\le 20\Rightarrow 3,3\le k_2\le 11,1\Rightarrow k_2=4,...,11:\text{ số giá trị } k_2 \text{ là 8} 
				\end{aligned}\right.$\\
				$\Rightarrow \text{Số vạch sáng}=12+8-4=16$.
			\end{itemize}
		}
	\end{ex}
	
	\begin{ex} %[3P5-3.2K][Loai7]
		Trong thí nghiệm Young về giao thoa ánh sáng đơn sắc, khoảng cách giữa $4$ vân sáng liên tiếp trên màn quan sát là $2,4\ mm$. Khoảng vân trên màn là
		\choice
		{$1,6\ mm$}	
		{$1,2\ mm$}
		{$0,6\ mm$}	
		{\True $0,8\ mm$}
		\loigiai{
			\begin{itemize}
				\item 4 vân sáng: $3i = 2,4 \Rightarrow i = 0,8\ mm$.
				\item 
			\end{itemize}
		}
	\end{ex}
	\begin{ex}%[3P2-8.2B][Loai1]
		Trên một sợi dây đàn hồi dài $1,05\ m$ với đầu A cố định, đầu B tự do đang có sóng dừng, tần số sóng là $100\ Hz$. Nếu không kể đầu B người ta thấy trên dây có $3$ bụng sóng. Tốc độ truyền sóng trên dây là
		\choice
		{$40\ m/s$}
		{$35\ m/s$}
		{$50\ m/s$}
		{\True $60\ m/s$}
		\loigiai{
			\begin{itemize}
				\item Ta có: $\ell = 3,5\dfrac{\lambda}{2} \Rightarrow \lambda = \dfrac{2\ell}{3,5} = \dfrac{2.1,05}{3,5} = 0,6\ m\Rightarrow v = \lambda.f = 0,6.100 = 60\ Hz.$
				\item 
				\item
			\end{itemize}
		}
	\end{ex}
	%\loai{Sự thay đổi của T}
	\begin{ex}%[3P2-8.2K][Loai1]
		Một sóng cơ truyền trên một sợi dây rất dài thì một điểm M trên sợi dây có vận tốc dao động biến thiên theo phương trình:  $v_M=20\pi \cos(10\pi t + \varphi )\ cm/s$. Giữ chặt một điểm trên dây sao cho trên dây hình thành sóng dừng và bỏ qua sự phản xạ qua lại nhiều lần. Khi đó bề rộng một bụng sóng có độ lớn là
		\choice
		{\True $8\ cm$}
		{$6\ cm$}
		{$16\ cm$}
		{$4\ cm$}
		\loigiai{
			\begin{itemize}
				\item Ta có phương trình vận tốc:
				$v = 20\pi.\cos(10\pi.t + \varphi)\Rightarrow A = 2\ cm $.
				\item Bề rộng của bụng sóng là $4A=8$ cm.
			\end{itemize}
		}
	\end{ex}
	\begin{ex}%[3P5-3.2K][Loai10]
		Trong thí nghiệm Young về giao thoa ánh sáng, khoảng cách giữa hai khe là $1,5\ mm$, khoảng cách từ hai khe đến màn là $3\ m$, người ta đo được khoảng cách giữa vân sáng bậc $2$ đến vân sáng bậc $5$ ở cùng phía với nhau so với vân sáng trung tâm là $3\ mm$. Bước sóng của ánh sáng dùng trong thí nghiệm là
		\choice
		{$0,2\ \mu m $}
		{$0,4\ \mu m $}
		{\True $0,5\ \mu m $}
		{$0,6\ \mu m $}
		\loigiai{
			\begin{itemize}
				\item Khoảng cách giữa vân sáng bậc 2 đến vân sáng bậc 5 cùng phía với nhau so với vân sáng trung tâm bằng:
				$5i-2i=3i=3\ mm\Rightarrow i=1mm\Rightarrow \lambda =\dfrac{ia}{D}=\dfrac{1. 1,5}{3}=0,5\ \mu m $.
				\item 
				\item 
			\end{itemize}
		}
	\end{ex}
	\begin{ex}
		\immini[thm]{Môi trường sống tự nhiên của dơi thường là rừng cây, và chúng bay xuyên qua các tán cây với tốc độ từ 10–15 km/h, tùy thuộc vào loài. Để bắt côn trùng, dơi cần có khả năng xác định chính xác vị trí của con mồi. Mặc dù thị giác của chúng tốt hơn so với câu nói "mù như dơi" (ám chỉ dơi không thấy gì), hệ thống định vị bằng tiếng vang của chúng sử dụng các xung âm thanh tần số rất cao $(50 - 100\ kHz)$, giúp chúng nhận biết chi tiết về thế giới ở các khoảng cách dưới 5 mét.}{\includegraphics[scale=0.04]{Images/115430}}
		Ở khoảng cách xa hơn, tiếng vang bị suy giảm quá nhiều khiến chúng không thể sử dụng. Dơi phát ra một tiếng "chíp" qua mũi. Xung âm thanh này thường kéo dài khoảng 3 ms. Khi âm thanh chạm vào các vật thể gần đó, nó sẽ phản xạ lại tai của dơi, và não của chúng có thể đo chính xác thời gian giữa phát ra xung và nhận được tiếng vang. Não của dơi cũng đã tiến hóa để tính toán khoảng cách đến vật phản xạ một cách tự nhiên. Biết tốc độ truyền âm trong không khí là 330 m/s và dơi nghe thấy âm thanh mình phát ra sau 0,02 giây thì con mồi cách dơi là
		\choice
		{3,8 m}
		{2,5 m}
		{\True 3,3 m}
		{4,5 m}
		\loigiai{
			\begin{itemize}
				\item $t=\dfrac{2L}{v}\Rightarrow L=\dfrac{vt}{2}=\dfrac{330.0,02}{2}=3,3\ m$.
			\end{itemize}
		}
		\haidong{15}
		
		
	\end{ex}
	\begin{ex} %[3P1Y5-1][Loai1]   
		Năng lượng vật dao động điều hòa
		\choice
		{tỉ lệ với biên độ dao động}
		{\True bằng thế năng của vật khi vật có li độ cực đại}
		{bằng thế năng của vật khi vật qua vị trí cân bằng}
		{bằng động năng của vật khi có li độ cực đại}
		\loigiai{
			Ta có: $W=\dfrac{1}{2}k.A^2$.
		}
		%\haidong{5}
	\end{ex}
	%\loai{Các yếu tố chu kì phụ thuộc}
	
	\begin{ex} %[3P1BF-1][Loai1]  
		Khi khối lượng của vật nặng là m thì chu kì dao động của con lắc đơn là $T$. Vậy khi tăng khối lượng của vật nặng lên $4$ lần thì chu kì dao động của con lắc đơn lúc này là
		\choice
		{$2T$}
		{\True $T$}
		{$4T$}
		{$\dfrac{T}{2}$}
		\loigiai{
			Chu kì của con lắc đơn được tính theo công thức: $T=2\pi \sqrt{{\dfrac{\ell}{g}}}$ không phụ thuộc vào khối lượng của vật nặng nên chu kì của con lắc đơn sau khi thay đổi khối lượng vật nặng vẫn là T.
		}	
		%\haidong{15}
	\end{ex}	
	%\loai{Biết chiều dài max, min tìm A, $\Delta \ell_0$ và năng lượng}
	
	
	\begin{ex}%[3P1K6-2][Loai1] 
		Một con lắc lò xo dao động điều hòa theo phương ngang. Trong quá trình dao động, chiều dài của lò xo biến thiên từ $22\ cm$ đến $30\ cm$. Biết lò xo của con lắc có độ cứng $100\ N/m$. Khi vật cách vị trí biên $3\ cm$ thì động năng của vật là
		\choice
		{$0,045\ J$}
		{$0,0375\ J$}
		{\True $0,075\ J$}
		{$0,035\ J$}
		\loigiai{
			\begin{itemize}
				\item Biên độ dao động của con lắc là: $A=\dfrac{30-22}{2}=4\ cm$
				\item Áp dụng định luật bảo toàn năng lượng ta có:
				\begin{align*}
					W_{\text{đ}}+W_t=\dfrac{kA^2}{2}\Rightarrow W_{\text{đ}}=\dfrac{kA^2}{2}-W_t=\dfrac{kA^2}{2}-\dfrac{kx^2}{2}=\dfrac{100.0,04^2}{2}-\dfrac{100.0,01^2}{2}=0,075\ J.
				\end{align*}
			\end{itemize}
		}
		%\haidong{15}
	\end{ex}
	\begin{ex} %[3P1-16.1B][Loai1]   
		Một con lắc đơn dao động điều hòa với biên độ góc $\alpha_0$. Lấy mốc thế năng ở vị trí cân bằng. Ở vị trí con lắc có động năng bằng thế năng thì li độ góc của nó bằng
		\choice
		{$\pm \dfrac{\alpha _0}{\sqrt{3}}$}
		{$\pm \dfrac{\alpha _0}{2}$}
		{\True $\pm \dfrac{\alpha _0}{\sqrt{2}}$}
		{$\pm \dfrac{\alpha _0}{3}$}
		\loigiai{Ở vị trí con lắc có động năng bằng thế năng thì li độ góc của nó bằng:
			\begin{align*}
				W_t=W_{\text{đ}}=\dfrac{W}{2}\Leftrightarrow \dfrac{mg\alpha^2}{2}=\dfrac{1}{2}\dfrac{mg\alpha_0^2}{2}\Rightarrow \alpha=\pm \dfrac{\alpha_0}{\sqrt{2}}
		\end{align*}}
		%\haidong{15}
	\end{ex}
	%\loai{Thế năng và động năng}
	\begin{ex} %[3P1-16.3K][Loai1]  
		Cho một con lắc đơn dao động điều hòa tại nơi có $g = 10\ m/s^2$. Biết rằng trong khoảng thời gian $12\ s$ thì nó thực hiện được $24$ dao động. Vận tốc cực đại của con lắc là $6\pi\ cm/s$. Lấy $\pi^2 = 10$. Giá trị góc lệch của con lắc so với phương thẳng đứng mà ở đó thế năng bằng $\dfrac{1}{8}$ động năng là
		\choice
		{$0,04\ rad$}
		{\True $0,08\ rad$}
		{$0,1\ rad$}
		{$0,12\ rad$}
		\loigiai{
			\begin{itemize}
				\item $T = \dfrac{12}{24} = 0,5\ s \Rightarrow \omega = \dfrac{2\pi}{0,5} = 4\pi \Rightarrow \ell = \dfrac{g}{\omega^2} = 0,0625\ m$ $\Rightarrow A = \dfrac{v_{\max}}{\omega} = 1,5\ cm$
				\item Ta có: $W_t = \dfrac{W}{8}_{\text{đ}} \Rightarrow W_t = \dfrac{W}{9} \Rightarrow |x| = \dfrac{A}{3} = 0,5\ cm = 0,005\ m$
				\item Mặt khác: $|x| = \ell|\alpha| \Rightarrow |\alpha| = \dfrac{|x|}{\ell} = \dfrac{0,005}{0,0625} = 0,08$ rad.
			\end{itemize}
		}
		%\haidong{15}
	\end{ex}
	

\subsubsection*{PHẦN 2. Thí sinh trả lời từ câu 1 đến câu 4. Trong mỗi ý a), b), c), d) ở mỗi câu, thí sinh chọn đúng hoặc sai.}
\setcounter{ex}{0}
\begin{ex}
	Một sợi dây đàn hồi dài $90 ~cm$ đang có sóng dừng với hai tần số liên tiếp là $30 ~Hz$ và $50 ~Hz$. Biết tốc độ truyền sóng là $40 ~cm/s$.
	\choiceTF[t]
		{Dây có một đầu cố định và một đầu tự do}
		{Tần số nhỏ nhất để có sóng dừng khi đó là $30 ~Hz$}
		{Bước sóng của sợi dây khi đang có sóng dừng với tần số nhỏ nhất là $15 ~cm$}
		{Khi đó trên dây có 10 nút sóng (kể cả hai nút ở hai đầu dây nếu có)}
	\loigiai{
		\begin{itemchoice}
		\itemch
		\itemch 
		\itemch 
		\itemch 
	\end{itemchoice}
	}
\end{ex}
\begin{ex}
	Sóng cơ là sự lan truyền dao động cơ trong môi trường vật chất đàn hồi. Khi sóng cơ truyền đi chỉ có pha dao động của các phần tử vật chất lan truyền còn các phần tử vật chất thì dao động xung quanh vị trí cân bằng cố định.
	\choiceTF[t]
	{Những phần tử của môi trường cách nhau một số nguyên lần bước sóng thì dao động cùng pha với nhau}
	{Hai phần tử của môi trường cách nhau một phần tư bước sóng trên cùng một phương truyền sóng thì dao động ngược pha nhau với nhau}
	{\True Những phần tử của môi trường trên cùng một hướng truyền sóng và cách nhau một số nguyên lần bước sóng thì dao động cùng pha với nhau}
	{Hai phần tử của môi trường trên cùng một phương truyền sóng cách nhau một nửa bước sóng thì dao động vuông pha với nhau}
	\loigiai{
		\begin{itemchoice}
			\itemch Sai, nhưng phần tử của môi trường cách nhau một số nguyên lần bước sóng dao động cùng pha khi cùng một phương truyền sóng
			\itemch Sai, hai phần tử của môi trường cách nhau một phần tư bước sóng trên cùng một phương truyền sóng thì dao động vuông pha nhau
			\itemch Đúng.
			\itemch Sai, hai phần tử của môi trường trên cùng một phương truyền sóng cách nhau một nửa bước sóng thì dao động ngược pha với nhau
		\end{itemchoice}
	}
\end{ex}
\begin{ex}
Một vật dao động điều hòa với phương trình li độ $x = A\cos(\omega t)$ có động năng và thế năng biến thiên liên tục theo thời gian, nhưng tổng cơ năng luôn được bảo toàn. 
	\choiceTF[t]
	{Đồ thị biểu diễn thế năng theo vận tốc là một đoạn thẳng đi qua gốc tọa độ}
	{Thế năng của chất điểm tại thời điểm $t$ có biểu thức là $W_{t}=\dfrac{W_0}{2}(1-\cos 2 \omega t)$}
	{\True Khi động năng của chất điểm giảm thì thế năng của nó tăng}
	{Động năng biến thiên vuông pha so với thế năng}
	\loigiai{
		\begin{itemchoice}
			\itemch Sai.
			\itemch Sai.
			\itemch Đúng.
			\itemch Sai.
		\end{itemchoice}
	}
	%\haidong{25}
\end{ex}
\begin{ex}
	Trong thí nghiệm giao thoa Young về giao thoa ánh sáng có bước sóng $\lambda$, khoảng cách từ màn tới hai khe là $1,5 ~m$, khoảng cách giữa hai khe là $a=1 ~mm$. Biết khoảng cách từ vân trung tâm tới vân sáng bậc 5 là $4,5 ~mm$.
	\choiceTF[t]
	{\True Khoảng vân là $i=0,9~mm$}
	{Bước sóng giao thoa là $\lambda=6000~nm$}
	{Điểm cách vân trung tâm $3,5 ~mm$ là vân sáng bậc 4}
	{Giá sử thay đổi nguồn sáng phát ra bước sóng $\lambda=0,4\ \mu m$ thì khoảng vân tăng 1,5 lần}
	
	\loigiai{
		\begin{itemchoice}
			\itemch 
			\itemch 
			\itemch 
			\itemch 
			
		\end{itemchoice}
	}
\end{ex}
\subsubsection*{PHẦN 3. Thí sinh trả lời từ câu 1 đến câu 6.}
\setcounter{ex}{0}
\begin{ex}
	Một chất điểm dao động điều hoà với phương trình vận tốc $v=10 \pi \cos (\pi t+\dfrac{\pi}{3})\ cm/s$ ($t$ tính bằng giây). Lấy $\pi^2=10$. Khi chất điểm có li độ là $5 \sqrt{3}\ cm$ thì tốc độ của chất điểm bằng bao nhiêu $cm/s$ (làm tròn kết quả đến chữ số hàng phần mười)? 
	\shortans[oly]{15,7}
	\loigiai{
	}
	%\haidong{30}
\end{ex}
\begin{ex}
	Một chất điểm dao động điều hòa theo phương trình $x=4 \cos \left(\pi t+\dfrac{\pi}{3}\right)$, với $x$ tính bằng $cm$ và $t$ tính bằng giây. Thời điểm vận tốc của chất điểm đạt giá trị $2 \pi\ cm/s$ lần thứ 5 là bao nhiêu giây (làm tròn kết quả đến chữ số hàng phần trăm)?
	\shortans[oly]{4,83} \loigiai{
		4,83
	}
	%\haidong{30}
\end{ex}
\begin{ex}
	Một chất điểm dao động điều hòa với phương trình li độ $x=10 \cos \left(4 \pi t+\dfrac{\pi}{6}\right)$
	(t tính bằng giây). Quãng đường chất điểm đi được kể từ thời điểm $t=0$ đến thời điểm $t=\dfrac{11}{6}\ s$ là bao nhiêu mét (làm tròn kết quả đến chữ số hàng phần trăm)?
	\shortans[oly]{1,49}
	\loigiai{
		$1,49 \ m$
	}
	%\haidong{30}
\end{ex}
\begin{ex}
	Một sóng hình sin truyền theo phương $Ox$ từ nguồn O với tần số $20 \ Hz$, có tốc độ truyền sóng nằm trong khoảng từ $0,6\ m/s$ đến $1 \ m/s$. Gọi A và B là hai điểm nằm trên $Ox$, ở cùng một phía so với O và cách nhau $10\ cm$. Hai phần tử môi trường tại A và B luôn dao động ngược pha với nhau. Tốc độ truyền sóng bằng bao nhiêu $m/s$? 
	\shortans[oly]{0,8} 
	\loigiai{
		Theo giả thiết hai phần tử môi trường tại A và B luôn dao động ngược pha với nhau
		$$
		\begin{aligned}
			& \Rightarrow \Delta \varphi=2 \pi \dfrac{d}{\lambda}=2 \pi \dfrac{d. f}{v}=(2 k+1) \pi \Rightarrow v=\dfrac{2 df}{2 k+1} \\
			& 0,6<\dfrac{2 df}{2 k+1}<1 \Rightarrow 1,5<k<2,8 \Rightarrow k=2 \Rightarrow v=0,8 \ m/s
		\end{aligned}
		$$
	}
\end{ex}
\begin{ex}
	AB là một sợi dây đàn hồi căng thẳng nằm ngang, M là một điểm trên AB cách A một khoảng bằng $12,5 \ cm$. Cho A dao động điều hòa và bắt đầu đi lên từ vị trí cân bằng. Biết bước sóng là $25 \ cm$ và tần số sóng là $5 \ Hz$. Kể từ khi A bắt đầu dao động, sau bao nhiêu giây thì M lên đến điểm cao nhất? 
	\shortans[oly]{0,15} 
	\loigiai{
		\begin{itemize}
			\item $v=\lambda f=25.5=125 \ cm/s$
			\item 			$t=\dfrac{d}{v}+\dfrac{T}{4}=\dfrac{d}{v}+\dfrac{1}{4 f}=\dfrac{12,5}{125}+\dfrac{1}{4.5}=0,15 \ s$
		\end{itemize}
	}
\end{ex}
\begin{ex}
	Lúc $t=0$ đầu O của dây cao su căng thẳng nằm ngang bắt đầu dao động đi lên với chu kì $2 \ s$, tạo thành sóng ngang lan truyền trên dây với tốc độ $2 \ cm/s$. Điểm M trên dây cách O một khoảng $1,4 \ cm$. Thời điểm đầu tiên để M đến điểm thấp nhất là bao nhiêu giây? 
	\shortans[oly]{2,2} 
	\loigiai{
		$t=\dfrac{d}{v}+\dfrac{3 \ T}{4}=\dfrac{1,4}{2}+\dfrac{3.2}{4}=2,2 \ s$
	}
\end{ex}